\documentclass[a4paper,11pt]{article}

% Pacotes básicos
\usepackage[utf8]{inputenc}
\usepackage[T1]{fontenc}
\usepackage[brazil]{babel}
\usepackage{geometry}
\usepackage{setspace}
\usepackage{titlesec}
\usepackage{amsmath, amssymb, amsfonts}
\usepackage{booktabs}
\usepackage{array}
\usepackage{graphicx}
\usepackage{float}
\usepackage{subcaption}
\usepackage[hidelinks]{hyperref}

\geometry{top=1.0cm, bottom=1.8cm, left=1.8cm, right=1.8cm}
\onehalfspacing

% Título
\title{\textbf{Relatório T1} \\ \large MC906}
\author{
    Adriano Ribeiro Franulovic Campos,
    Andrei Quaresma Pinto,
    Matheus Ferracciú Scatolin,\\
    Pietro Fernandes Magaldi,
    Tiago Perrupato Antunes
}
\date{}

\begin{document}

\maketitle

\section{Modelagem do problema}

O ambiente foi implementado como um MDP em grade. A malha padrão possui dimensão $5\times5$ e as células usam a codificação:
\begin{itemize}
    \item \texttt{S}: estado inicial,
    \item \texttt{G}: objetivo terminal,
    \item \texttt{T}: armadilha terminal,
    \item \texttt{\#}: parede (estado bloqueado),
    \item \texttt{.}: estado livre.
\end{itemize}

No layout padrão, existem 25 posições na grade e 5 paredes, portanto o conjunto de estados utilizáveis tem cardinalidade $|S|=20$. Os estados terminais são dois (\texttt{G} e \texttt{T}). Na UI também foram disponibilizados layouts maiores (\texttt{OPEN 7x7}, \texttt{MAZE 9x9} e \texttt{GAUNTLET 6x10}), preservando a mesma semântica de células e recompensas.

\textbf{Estados ($S$).} Cada estado é um par de coordenadas $s=(r,c)$ para toda célula não bloqueada.

\textbf{Ações ($A$).} O conjunto de ações é discreto e fixo:
\[
A=\{\text{UP},\text{DOWN},\text{LEFT},\text{RIGHT}\}.
\]

\textbf{Dinâmica ($T$).} A função de transição é fornecida por:
\begin{itemize}
    \item modo determinístico (padrão): probabilidade 1 para o movimento escolhido;
    \item modo estocástico: a ação desejada ocorre com probabilidade \texttt{wind\_prob} (padrão 0.8), e as duas ações perpendiculares dividem igualmente a massa restante.
\end{itemize}

Caso o agente tente sair da grade ou bata na parede, ele permanece no mesmo estado. Em estados terminais, a transição é absorvente (retorna ao próprio estado com probabilidade 1 e recompensa imediata 0).

\textbf{Recompensas e utilidade ($U$).} A recompensa imediata é definida por:
\[
R(s,a,s')=
\begin{cases}
+1.0, & s'\in G \\
-1.0, & s'\in T \\
-0.04, & \text{caso contrário}
\end{cases}
\]

O termo de utilidade de longo prazo é o retorno descontado esperado:
\[
U(s)=\mathbb{E}\left[\sum_{k=0}^{\infty}\gamma^k r_{t+k+1}\mid s_t=s\right].
\]

A penalidade por passo ($-0.04$) força uma política que evita trajetórias longas e ajuda a diferenciar caminhos de mesmo destino.

\section{Equações de Bellman}

Foi implementada Value Iteration com atualização síncrona. Para cada estado não terminal, calcula-se:
\[
V_{k+1}(s)=\max_{a\in A}\sum_{s'} T(s,a,s')\left[R(s,a,s')+\gamma V_k(s')\right].
\]

Na implementação, para cada ação é computado um valor intermediário
\[
Q_k(s,a)=\sum_{s'} T(s,a,s')\left[R(s,a,s')+\gamma V_k(s')\right],
\]
e depois selecionado o máximo em relação a $a$. Para estados terminais, o código fixa $V(s)=0$ (a recompensa de atingir terminal já foi contabilizada na transição anterior).

\textbf{Critério de convergência.} Em cada iteração:
\[
\Delta_k=\max_{s\in S}|V_{k+1}(s)-V_k(s)|.
\]
O algoritmo para quando $\Delta_k<\theta$ (padrão $\theta=10^{-6}$) ou ao atingir \texttt{max\_iterations} (100000).

\textbf{Extração de política.} Após convergir, a política gulosa é extraída por:
\[
\pi^*(s)=\arg\max_{a\in A}\sum_{s'} T(s,a,s')\left[R(s,a,s')+\gamma V(s')\right].
\]

\section{Q-learning}

Foi utilizado Q-learning com aprendizado tabular. A tabela é inicializada como:
\[
Q(s,a)=0\quad \forall (s,a),
\]
com dimensão $(|S|,4)$, onde as quatro colunas correspondem às ações.

\subsection*{Regra de atualização}

Para cada transição $(s,a,r,s')$, aplica-se:
\[
Q(s,a)\leftarrow Q(s,a)+\alpha\Big(r+\gamma\max_{a'}Q(s',a')-Q(s,a)\Big).
\]

O episódio inicia em \texttt{S}, termina ao atingir estado terminal ou \texttt{max\_steps} (padrão 200), e registra retorno acumulado e número de passos.
\section{Double Q-learning}
Além do Q-learning tradicional, foi implementada a variante Double Q-learning, que mantém duas tabelas $Q_A$ e $Q_B$ para reduzir o viés de superestimação. A cada passo, uma das tabelas é escolhida aleatoriamente para atualização, usando a outra tabela para estimar o valor futuro:
\[
\begin{cases}
Q_A(s,a)\leftarrow Q_A(s,a)+\alpha\Big(r+\gamma Q_B(s',\arg\max_{a'}Q_A(s',a'))-Q_A(s,a)\Big), & \text{com prob. } 0.5 \\
Q_B(s,a)\leftarrow Q_B(s,a)+\alpha\Big(r+\gamma Q_A(s',\arg\max_{a'}Q_B(s',a'))-Q_B(s,a)\Big), & \text{com prob. } 0.5
\end{cases}
\]
Onde a política gulosa é extraída por:
\[
\pi(s)=\arg\max_a \left(Q_A(s,a)+Q_B(s,a)\right).
\]

O episódio e os critérios de término são os mesmos do Q-learning tradicional.

\section{Exploração adotada}

Para ambas as estratégias (Q-learning e Double Q-learning), a política de comportamento é $\varepsilon$-gulosa:
\begin{itemize}
    \item com probabilidade $\varepsilon$, escolhe ação aleatória uniforme;
    \item caso contrário, escolhe $\arg\max_a Q(s,a)$.
\end{itemize}

Ao final de cada episódio:
\[
\varepsilon\leftarrow\max(\varepsilon_{\min},\varepsilon\cdot \text{decay}).
\]

Valores padrão do experimento base: $\alpha=0.1$, $\gamma=0.99$, $\varepsilon_0=1.0$, \texttt{decay}=0.995 e $\varepsilon_{\min}=0.05$.

Na interface interativa, o treinamento de Q-learning usa por padrão 500 episódios e exibe no HUD o progresso por episódio, retorno acumulado, número de passos, valor de $\varepsilon$ e velocidade de simulação.

\subsection*{Impacto esperado dos parâmetros}

\begin{itemize}
    \item $\alpha$ controla velocidade de adaptação: maior $\alpha$ aprende rápido, porém com maior variância; menor $\alpha$ estabiliza, mas torna o processo lento.
    \item $\gamma$ controla horizonte: valores altos favorecem retornos futuros e tendem a políticas de longo prazo; valores baixos privilegiam ganho imediato.
    \item $\varepsilon$ e \texttt{decay} regulam exploração/explotação: muita exploração reduz risco de ótimo local, mas piora retorno inicial; pouca exploração acelera desempenho inicial, porém pode cristalizar política subótima.
\end{itemize}

\section{Visualização}

Além dos plots estáticos, foi implementada uma UI arcade que permite acompanhar o treinamento em tempo real e inspecionar qualitativamente a política aprendida.

Na tela de menu, é possível configurar diferentes experimentos combinando: (i) algoritmo (\textit{Q-learning}, (\textit{Double Q-learning} ou \textit{Value Iteration}), (ii) mapa (\texttt{DEFAULT 5x5}, \texttt{OPEN 7x7}, \texttt{MAZE 9x9}, \texttt{GAUNTLET 6x10}) e (iii) vento estocástico ligado/desligado. Durante a execução, a UI oferece controles de pausa, passo-a-passo e velocidade, além de sobreposições de valor ($V$), política ($P$) e mapa de visitas ($H$). Ao final, é possível reproduzir a trajetória gulosa e gravar GIFs para documentação dos resultados.

\section{Avaliação Experimental e Discussão}

Os experimentos realizados avaliaram o desempenho e a convergência dos algoritmos \textit{Value Iteration} (VI), \textit{Q-learning} e (\textit{Double Q-learning} no ambiente \textit{GridWorld} proposto. A análise conjunta dos resultados revela grande consistência entre a abordagem \emph{model-based} e a \emph{model-free}, com ambas alcançando uma política de alta qualidade, mas exibindo dinâmicas de aprendizado e \textit{trade-offs} distintos ao longo do processo. A UI foi usada como apoio para validação qualitativa das trajetórias e da evolução das políticas em diferentes mapas.

\textbf{Convergência e Estabilidade.} O algoritmo \textit{Value Iteration} apresentou uma convergência notavelmente rápida e com decréscimo monotônico do erro (observável pela curva de $\max|\Delta V|$), estabilizando completamente em cerca de 10 iterações. Esse comportamento é esperado, dado que o algoritmo possui acesso explícito e integral às funções de transição e recompensa do modelo, além de operar sobre um espaço de estados de dimensionalidade reduzida. Em contrapartida, a curva de aprendizado do \textit{Q-learning} reflete sua natureza amostral e exploratória. Durante os episódios iniciais, o algoritmo registra alta variância e retornos fortemente negativos (com quedas próximas a $-8$), decorrentes das frequentes visitas à armadilha geradas pela exploração aleatória. No entanto, o aprendizado demonstra-se eficaz: após uma fase transitória que varia entre 100 e 250 episódios, o retorno estabiliza-se em um patamar positivo constante, equiparando-se ao ótimo obtido analiticamente pelo VI.

\textbf{Mapas de Valor e Políticas Ótimas.} Ao final do treinamento, tanto o planejamento do VI quanto o aprendizado do \textit{Q-learning} produziram trajetórias gulosas idênticas. Ambos os agentes aprendem a se deslocar do estado inicial até o objetivo de forma ótima (alcançando o retorno de $0.760$), contornando os obstáculos e mantendo uma distância segura da armadilha. A análise dos mapas espaciais da função de valor ($V(s)$) confirma esse entendimento estratégico: há um claro gradiente crescente de utilidade direcionado à rota do objetivo (\texttt{G}), enquanto os estados adjacentes à armadilha (\texttt{T}) ou de trânsito custoso recebem valorações inferiores. Essa paridade espacial nos mapas de valor valida que o \textit{Q-learning} obteve uma representação interna do ambiente muito próxima daquela computada exaustivamente pelo VI.

\textbf{Impacto dos Hiperparâmetros.} A variação na estratégia de exploração do \textit{Q-learning} elucida o clássico dilema \textit{exploration vs. exploitation}. Configurações com alta taxa de exploração inicial (como $\varepsilon = 1.0$) garantem uma varredura abrangente do espaço de estados, reduzindo o risco de o agente cristalizar prematuramente uma política subótima (miopia estratégica). O custo inerente dessa abordagem é um pior desempenho inicial, evidenciado pelas frequentes quedas na recompensa. Em contrapartida, variações no fator de desconto $\gamma$ (0.5, 0.9 e 0.99) demonstraram impactar primordialmente a velocidade com que o algoritmo cruza a fase transitória de estabilização, embora, neste ambiente específico, todas as configurações testadas tenham convergido para patamares de retorno médio muito semelhantes ao final dos 1000 episódios.

\textbf{Análise com diferentes sementes.} Quanto à análise em diferentes sementes aleatórias de grids, temos que a agregação por múltiplas sementes confirma que as flutuações de desempenho — marcantes na fase transitória — são inerentes à estocasticidade do aprendizado \emph{model-free}. No entanto, reafirma a robustez assintótica do algoritmo: independentemente do caminho de exploração escolhido (rápido e direto ou demorado e abrangente), o \textit{Q-learning} é capaz de superar a ignorância inicial e encontrar a trajetória ótima de forma confiável.

\textbf{Trade-offs e Limitações.} Os resultados materializam os principais \textit{trade-offs} da literatura: métodos \emph{model-based} oferecem estabilidade e rapidez algorítmica, porém ao custo restritivo de exigir um modelo perfeito (\textit{prior knowledge}) do domínio. Métodos \emph{model-free}, por dispensarem essa premissa, operam com maior flexibilidade, mas demandam um custo amostral significativamente mais alto no início da execução. 

\section{Plots}

\begin{figure}[b]
    \centering
    \begin{subfigure}[b]{0.5\textwidth}
        \centering
        \includegraphics[width=\textwidth]{ql/learning_curve.png}
        \caption{Curva de aprendizado}
        \label{fig:plot1}
    \end{subfigure}
    \hfill 
    \begin{subfigure}[b]{0.35\textwidth}
        \centering
        \includegraphics[width=\textwidth]{ql/value_map.png}
        \caption{Mapa de valores}
        \label{fig:plot2}
    \end{subfigure}
    \hfill 
    \begin{subfigure}[b]{0.35\textwidth}
        \centering
        \includegraphics[width=\textwidth]{ql/policy.png}
        \caption{Política aprendida}
        \label{fig:plot2}
    \end{subfigure}
    \hfill 
    \begin{subfigure}[b]{0.35\textwidth}
        \centering
        \includegraphics[width=\textwidth]{ql/trajectory.png}
        \caption{Trajetória aprendida}
        \label{fig:plot2}
    \end{subfigure}
    \caption{Plots obtidos para algoritmo de Q-learning.}
    \label{fig:conjunto Q-learning}
\end{figure}

\begin{figure}
    \centering
    \begin{subfigure}[b]{0.5\textwidth}
        \centering
        \includegraphics[width=\textwidth]{vi/convergence.png}
        \caption{Curva de convergência}
        \label{fig:plot1}
    \end{subfigure}
    \hfill 
    \begin{subfigure}[b]{0.35\textwidth}
        \centering
        \includegraphics[width=\textwidth]{vi/value_map.png}
        \caption{Mapa de valores}
        \label{fig:plot2}
    \end{subfigure}
    \hfill 
    \begin{subfigure}[b]{0.35\textwidth}
        \centering
        \includegraphics[width=\textwidth]{vi/policy.png}
        \caption{Política aprendida}
        \label{fig:plot2}
    \end{subfigure}
    \hfill 
    \begin{subfigure}[b]{0.35\textwidth}
        \centering
        \includegraphics[width=\textwidth]{vi/trajectory.png}
        \caption{Trajetória aprendida}
        \label{fig:plot2}
    \end{subfigure}
    \caption{Plots obtidos para algoritmo de Value Iteration.}
    \label{fig:conjunto Value Iteration}
\end{figure}

\begin{figure}
    \centering
    \begin{subfigure}[b]{0.44\textwidth}
        \centering
        \includegraphics[width=\textwidth]{exploration_comparison.png}
        \label{fig:plot1}
    \end{subfigure}
    \hfill 
    \begin{subfigure}[b]{0.44\textwidth}
        \centering
        \includegraphics[width=\textwidth]{gamma_comparison.png}
        \label{fig:plot2}
    \end{subfigure}
    \hfill 
    \begin{subfigure}[b]{0.44\textwidth}
        \centering
        \includegraphics[width=\textwidth]{vi_vs_ql_rewards.png}
        \label{fig:plot2}
    \end{subfigure}
    \hfill 
    \begin{subfigure}[b]{0.44\textwidth}
        \centering
        \includegraphics[width=\textwidth]{exploration.jpeg}
        \label{fig:seeds}
    \end{subfigure}
    \caption{Comparação entre algoritmos implementados.}
    \label{fig:comparacao}
\end{figure}

\end{document}
